\chapter{Het zenuwstelsel beschermen}\label{ch:g8:protecting-the-nervous-system}

Twee hoofdstukken hebben laten zien wat het \omterm{def:g5:brain-in-command:system}{zenuwstelsel} is: bedrading
waarvan de koppelingen leren, beslissen --- en te misleiden zijn
(\cref{rem:g8:nervous-communication:vulnerable}). De afsluiting van dit
jaar zet die anatomie om in bescherming. Wat doen slaaptekort, lawaai,
schermen en de geestbeïnvloedende stoffen eigenlijk, koppeling voor
koppeling, met de machine die jij bent --- en hoe ziet haar beschermen er
in de praktijk uit, op je veertiende?

\section{Onderhoud: slaap, en de grenzen van de zintuigen}

\begin{proposition}[Slaap is onderhoud van synapsen]\label{prop:g8:protecting-the-nervous-system:sleep}
Het opruimen en herstellen uit
\cref{prop:g5:brain-in-command:protect} laat zich op koppelingsniveau
lezen: in de \omterm{prop:g1:healthy-habits:needs}{slaap} worden de synapsen die overdag versterkt zijn
vastgezet --- ``oefening maakt blijvend''
(\cref{ex:g8:nervous-communication:juggler}) gebeurt grotendeels 's nachts
--- en ruimt het huishouden van de \omterm{def:g5:brain-in-command:system}{hersenen} de chemische slijtage van de
dag op. Korte nachten laten het leren van gisteren halfgearchiveerd en het
afwegen van morgen traag: de puberende \omterm{def:g5:brain-in-command:system}{hersenen}, midden in hun
herbedrading (\cref{rem:g8:puberty:moods}) en met een klok die naar achteren
schuift, hebben hun uren juist het hardst nodig als de school het vroegst
begint.
\end{proposition}
\admitted

\begin{example}[De leerparadox]\label{ex:g8:protecting-the-nervous-system:revision}
Twee leerlingen voor een toets: de een leert tot middernacht en gaat door
tot twee uur; de ander leert tot tien uur en gaat slapen. Meestal wint de
tweede --- de stof stak hun koppelingen over \emph{en werd vastgezet}; de
extra uren van de eerste kostten precies het archiveren dat het leren laat
beklijven. ``Er een nachtje over slapen'' is techniek van synapsen.
\end{example}

\begin{proposition}[Hard is een dosis]\label{prop:g8:protecting-the-nervous-system:noise}
De zintuigcellen van het \omterm{def:g1:the-five-senses:sense}{oor} --- de haarfijne ontvangers die geluid in
\omterm{def:g8:nervous-communication:neuron}{zenuwimpulsen} omzetten --- groeien niet terug
(\cref{prop:g1:the-five-senses:fragile} zei het voorzichtig; hier is het
mechanisme). Hard geluid doodt ze per dosis: volume maal blootstelling.
Concerten en oordopjes op vol volume geven uur na uur een kapitaal aan
ontvangers uit dat geen enkele nacht \omterm{prop:g1:healthy-habits:needs}{slaap} teruggeeft --- het fluiten na
een luide avond (\cref{ex:g1:the-five-senses:loud}) is de
overdosiswaarschuwing. De regel is een budget: zachter, korter, met pauzes
--- en oordoppen waar de versterkers heersen.
\end{proposition}
\admitted

\section{De stoffen die de koppelingen misleiden}

\begin{proposition}[Geestbeïnvloedende stoffen grijpen aan bij synapsen]\label{prop:g8:protecting-the-nervous-system:substances}
Elke stof die stemming, waarneming of alertheid verandert, werkt in bij de
spleten uit \cref{def:g8:nervous-communication:synapse} --- door een
boodschapper na te bootsen, er een te blokkeren, of afgifte af te dwingen:
\begin{itemize}
  \item \emph{alcohol} dempt de koppelingen van het hele stelsel: de lus
        vertraagt van begin tot eind --- oordeel, \omterm{def:g5:brain-in-command:reaction}{reactietijd} (de meters
        uit \cref{ex:g5:brain-in-command:driver} worden langer), evenwicht,
        geheugen; de nog in aanbouw zijnde koppelingen van puberende
        \omterm{def:g5:brain-in-command:system}{hersenen} krijgen de demping het diepst binnen;
  \item de \emph{nicotine} uit tabak bootst een natuurlijke boodschapper
        na en herbedraadt bij herhaling de beloningskoppelingen richting
        \emph{\omterm{prop:g5:healthy-choices:substances}{verslaving}} (de val uit
        \cref{prop:g5:healthy-choices:substances}, nu met haar mechanisme:
        de synapsen stellen zich in op de dosis --- en protesteren zonder);
  \item \emph{cannabis} bootst een andere familie boodschappers na en
        vertroebelt juist de koppelingen van geheugen, motivatie en timing
        --- opnieuw het duurst in \omterm{def:g5:brain-in-command:system}{hersenen} die in aanbouw zijn;
  \item hoe harder de stof, hoe grover het forceren van de koppelingen ---
        en hoe sneller de beloningsherbedrading die de motor van de
        \omterm{prop:g5:healthy-choices:substances}{verslaving} is.
\end{itemize}
\end{proposition}
\admitted

\begin{remark}[Verslaving, mechanisch bekeken]\label{rem:g8:protecting-the-nervous-system:addiction}
De val uit \cref{rem:g5:healthy-choices:never} heeft nu een
bedradingsschema: de beloningskoppelingen die herhaaldelijk gedoseerd
worden, \emph{passen zich aan} --- receptoren teruggesnoeid, drempels
verhoogd --- tot de stof nodig is om je normaal te voelen, en haar
afwezigheid een protest van het stelsel is. Die \omterm{def:g4:adapted-to-their-world:adaptation}{aanpassing} is de
scheefheid van makkelijke-ingang-moeilijke-uitgang, uit synapsen
opgebouwd: beginnen is één beslissing; stoppen gaat tegen herbedrade
koppelingen in. De nog goed instelbare puberende \omterm{def:g5:brain-in-command:system}{hersenen} herbedraden het
snelst --- en daarom vindt elke statistiek dat wie vroeg begint het diepst
vastzit, en daarom blijft nooit-beginnen de goedkope zet.
\end{remark}

\begin{example}[De schermlus]\label{ex:g8:protecting-the-nervous-system:screens}
Geen molecuul, dezelfde koppelingen: tijdlijnen en spelletjes zijn ---
met opzet --- afgestemd om de beloningsbedrading te kietelen met
onvoorspelbare kleine prijzen, het schema dat koppelingen het hardst
traint. Het gevolg is een milde neef van de val: checken dat zich tegen
stoppen verzet, avonden die verdampen, \omterm{prop:g1:healthy-habits:needs}{slaap} die naar achteren geduwd
wordt (de uren uit \cref{prop:g5:healthy-choices:screens}, plus de kosten
uit \cref{prop:g8:protecting-the-nervous-system:sleep}). Dezelfde
verdediging als altijd: beslis de uren vooraf, en houd het laatste uur
schermvrij.
\end{example}

\section{In de praktijk: je eigen bescherming draaien}

\begin{method}[De gebruikershandleiding van het zenuwstelsel]\label{met:g8:protecting-the-nervous-system:manual}
De zorglijst uit \cref{met:g5:healthy-choices:run}, opnieuw uitgegeven met
mechanismen erbij:
\begin{enumerate}
  \item \emph{\omterm{prop:g1:healthy-habits:needs}{slaap} de uren} --- het vastzetten en het huishouden lopen
        's nachts; bescherm het laatste uur tegen schermen;
  \item \emph{begroot het lawaai} --- kapitaal aan ontvangers groeit niet
        terug: volume omlaag, pauzes ertussen, oordoppen waar het buldert;
  \item \emph{helmen} waar vallen snel gaat
        (\cref{ex:g5:brain-in-command:helmet}) --- bedrading geneest niet
        zoals \omterm{def:g3:bones-and-muscles:skeleton}{bot};
  \item \emph{geen koppelingsmisleiders} voor \omterm{def:g5:brain-in-command:system}{hersenen} die in aanbouw zijn
        --- en de ``nee'' rustig van tevoren beslist
        (de zin op zak uit \cref{ex:g5:healthy-choices:answers}, nu met
        haar volledige onderbouwing);
  \item \emph{gebruik de machine} --- leren, sport, muziek: koppelingen
        worden sterker van verkeer, en een getrainde lus is het bezit dat
        elke andere gewoonte beschermt.
\end{enumerate}
\end{method}

\begin{example}[Hulp bestaat, en werkt]\label{ex:g8:protecting-the-nervous-system:help}
Voor wie al verstrikt zit --- in een eigen gewoonte of die van een vriend,
met stoffen of met schermen: de adressen uit
\cref{ex:g8:contraception:help} gelden ook hier. Schoolverpleegkundigen en
artsen behandelen \omterm{prop:g5:healthy-choices:substances}{verslaving} als wat dit hoofdstuk laat zien dat ze is ---
een bedradingsprobleem, geen oordeel over je karakter --- en herbedrade
koppelingen kunnen zich opnieuw aanpassen: moeilijker dan nooit beginnen,
en verre van hopeloos. Het slechtste plan is dat van het schoolplein:
zwijgen.
\end{example}

\section{Oefeningen}

\begin{exercise}[$\star$]\label{exo:g8:protecting-the-nervous-system:1}
Wat doet \omterm{prop:g1:healthy-habits:needs}{slaap} bij de koppelingen, en wie heeft de uren het hardst nodig?
\end{exercise}

\begin{exercise}[$\star$]\label{exo:g8:protecting-the-nervous-system:2}
Waarom is hard geluid een budgetkwestie? Wat wordt er uitgegeven, en wat is
de overdosiswaarschuwing?
\end{exercise}

\begin{exercise}[$\star$]\label{exo:g8:protecting-the-nervous-system:3}
Waar grijpen alle geestbeïnvloedende stoffen aan, en op welke drie
manieren?
\end{exercise}

\begin{exercise}[$\star$]\label{exo:g8:protecting-the-nervous-system:4}
Wat doet alcohol met de lus --- noem drie stations.
\end{exercise}

\begin{exercise}[$\star$]\label{exo:g8:protecting-the-nervous-system:5}
Beschrijf het mechanisme van \omterm{prop:g5:healthy-choices:substances}{verslaving} in twee zinnen.
\end{exercise}

\begin{exercise}[$\star$]\label{exo:g8:protecting-the-nervous-system:6}
Waarom hoort de schermlus in dit hoofdstuk thuis, molecuulvrij als ze is?
\end{exercise}

\begin{exercise}[$\star\star$]\label{exo:g8:protecting-the-nervous-system:7}
Verklaar de leerparadox met
\cref{prop:g8:protecting-the-nervous-system:sleep}.
\end{exercise}

\begin{exercise}[$\star\star$]\label{exo:g8:protecting-the-nervous-system:8}
Waarom kosten dezelfde doses puberende \omterm{def:g5:brain-in-command:system}{hersenen} meer dan die van een
volwassene? Twee mechanismen
(\cref{rem:g8:puberty:moods},
\cref{rem:g8:protecting-the-nervous-system:addiction}).
\end{exercise}

\begin{exercise}[$\star\star$]\label{exo:g8:protecting-the-nervous-system:9}
Verbind de meters uit \cref{ex:g5:brain-in-command:driver} met de demping
door alcohol: wat wordt er langer, en waarom is de wet op rijden onder
invloed een wet over \omterm{def:g5:brain-in-command:reaction}{reactietijd}?
\end{exercise}

\begin{exercise}[$\star\star$]\label{exo:g8:protecting-the-nervous-system:10}
Onderbouw drie punten van de gebruikershandleiding opnieuw met hun
mechanismen op koppelingsniveau.
\end{exercise}

\begin{exercise}[$\star\star$]\label{exo:g8:protecting-the-nervous-system:11}
``Een bedradingsprobleem, geen oordeel over je karakter.'' Wat verandert
die herformulering aan het zoeken van hulp --- en aan het aanbieden ervan
aan een vriend?
\end{exercise}

\begin{exercise}[$\star\star\star$]\label{exo:g8:protecting-the-nervous-system:12}
Schrijf het slotbetoog van het jaar: leid uit de instelbare spleten van
\cref{def:g8:nervous-communication:synapse} zowel de grootste kracht van de
\omterm{def:g5:brain-in-command:system}{hersenen} als hun kenmerkende kwetsbaarheden af --- leren en \omterm{prop:g5:healthy-choices:substances}{verslaving} als
één mechanisme dat je twee keer leest.
\end{exercise}

\section{Probleem: De preventieavond}

\begin{problem}[{Weekendprobleem --- de ouder-en-leerlingavond van de
school, met mechanismen uitgeschreven}]\label{pb:g8:protecting-the-nervous-system:1}
De klas bereidt de preventieavond van de school voor: vier stands ---
\omterm{prop:g1:healthy-habits:needs}{slaap}, geluid, stoffen, schermen --- elk met een demonstratie en een
mechanismekaart.

\medskip\noindent\textbf{Deel I --- De kaarten van de
stands.}
\begin{enumerate}
  \item Schrijf de slaapkaart: de demonstratie is de liniaaltest
        (\cref{met:g5:brain-in-command:ruler}) bij twee vrijwilligers,
        uitgerust tegenover met een korte nacht. Voorspel de uitkomsten en
        geef het mechanisme bij de koppelingen.
  \item Schrijf de geluidskaart: een decibelmeter, een paar oordoppen, en
        de zin ``je kunt ze niet terug laten groeien''. Vul het mechanisme
        in.
  \item Schrijf de stoffenkaart: één tekening van een \omterm{def:g8:nervous-communication:synapse}{synaps}, drie pijlen
        met de labels nabootsen, blokkeren en forceren. Wijs alcohol,
        nicotine en een gif uit
        \cref{exo:g8:nervous-communication:11} aan hun pijl toe.
  \item Schrijf de schermenkaart: geen molecuul --- wat deelt de
        demonstratie met een telefoon vol meldingen dan met de stand over
        stoffen?
\end{enumerate}

\medskip\noindent\textbf{Deel II --- De vragen uit de
zaal.}
\begin{enumerate}[resume]
  \item Een ouder: ``wij dronken op die leeftijd ook en het is goed met
        ons gekomen.'' Antwoord met de scheefheid van de bouwplaats ---
        wat is er anders op je veertiende?
  \item Een leerling: ``één sigaret kan niets herbedraden.'' Antwoord met
        \cref{rem:g8:protecting-the-nervous-system:addiction} en
        \cref{rem:g5:healthy-choices:never} --- wat is de werkelijke prijs
        van de eerste dosis?
  \item Een ouder: ``zijn schermen niet gewoon de nieuwe televisie?''
        Onderscheid de schema's: wat maakt onvoorspelbare kleine prijzen
        de hardere trainer?
  \item Een leerling, zachtjes: ``en als iemand al niet meer kan
        stoppen?'' Geef het antwoord van
        \cref{ex:g8:protecting-the-nervous-system:help}, met adressen erbij.
\end{enumerate}

\medskip\noindent\textbf{Deel III --- De slottoespraak.}
\begin{enumerate}[resume]
  \item De toespraak opent met de jongleur en sluit met \omterm{prop:g5:healthy-choices:substances}{verslaving} ---
        ``dezelfde instelbaarheid, twee keer gelezen''. Schrijf dat betoog
        in vier zinnen.
  \item Daarna krijgt elke leerling de zin op zak uit
        \cref{ex:g5:healthy-choices:answers} mee. Leg, mechanisme
        inbegrepen, uit waarom die zin \emph{nu} bedacht moet worden.
  \item De poster van de avond toont de vijf punten van de
        gebruikershandleiding. Zet ze op volgorde voor de week van een
        veertienjarige, het meest beschermende eerst, en verdedig je
        eerste keuze.
  \item Sluit de avond af in één zin: wat wordt er beschermd, en waarom is
        veertien het midden van het beslissende decennium?
\end{enumerate}
\end{problem}
